@@M0@@ to @@M1@@. This homomorphism defines a closed immersion @@M2@@, regular by 1.8 ii).

By 1.8 iii), @@M3@@ is perfect as a graded @@M4@@-Module. In particular, @@M5@@ is finitely presented as a graded Module. Let therefore:
@@M6@@
be a finite graded presentation of @@M7@@ over @@M8@@. On @@M9@@ we obtain the exact sequence:
@@M10@@
where @@M11@@, @@M12@@, so these are locally free of finite type over @@M13@@. We can therefore write, for every @@M14@@, the commutative diagram:
@@M15@@
where @@M16@@ is the structural morphism: @@M17@@, and @@M18@@, @@M19@@ are the degree @@M20@@ components of @@M21@@, @@M22@@. The upper row is exact by hypothesis; the first two vertical arrows are isomorphisms by the classical results for the direct images of @@M23@@, @@M24@@ being a projective bundle (EGA III, 2.1.15); finally, the last row is exact when @@M25@@ is sufficiently large, since the kernel of @@M26@@ is a perfect @@M27@@-Module, and we can apply the lemma recalled above, which shows that this kernel, tensored with @@M28@@ for @@M29@@ sufficiently large, is @@M30@@-acyclic; likewise for the kernel of @@M31@@. It follows that for @@M32@@ sufficiently large, @@M33@@ is an isomorphism. Since @@M34@@, we therefore have
